\section*{Capítulo \ref{ch:g11:reaction-energy} --- La energía de las reacciones: combustión y energías de enlace}

\begin{solution}{exo:g11:reaction-energy:1}
La madera que \omterm{def:g4:what-a-fire-needs:burning}{arde}: \omterm{def:g11:reaction-energy:exothermic}{exotérmica}. El hielo que se funde: \omterm{def:g11:reaction-energy:exothermic}{endotérmica} (toma
calor de su entorno, aunque no es una \omterm{def:g7:chemical-reactions:reaction}{reacción química}). La bolsa de
frío: \omterm{def:g11:reaction-energy:exothermic}{endotérmica}. El calentador de manos: \omterm{def:g11:reaction-energy:exothermic}{exotérmica}. Fabricar azúcar
con la luz del Sol: \omterm{def:g11:reaction-energy:exothermic}{endotérmica} (la energía procede de la luz).
\end{solution}

\begin{solution}{exo:g11:reaction-energy:2}
$Q = 3.0 \times 890.7 = \qty{2.7e3}{kJ}$, unos \qty{2.7}{MJ}.
\end{solution}

\begin{solution}{exo:g11:reaction-energy:3}
Los \omterm{def:g7:chemical-reactions:reactant}{reactivos}. La diferencia se libera al entorno, sobre todo en forma
de calor.
\end{solution}

\begin{solution}{exo:g11:reaction-energy:4}
$n = 100 / 46.0 = \qty{2.17}{mol}$; $Q = 2.17 \times 1367.6 =
\qty{2.97e3}{kJ}$, unos \qty{3.0}{MJ}.
\end{solution}

\begin{solution}{exo:g11:reaction-energy:5}
La energía necesaria para romper un \omterm{def:g10:the-mole:amount}{mol} de esos enlaces, estando las
\omterm{def:g7:atoms-and-molecules:molecule}{moléculas} y los \omterm{def:g7:atoms-and-molecules:atom}{átomos} en estado gaseoso. El par compartido mantiene
unidos los dos \omterm{def:g7:atoms-and-molecules:atom}{átomos}: separarlos va en contra de esa atracción, lo que
requiere energía.
\end{solution}

\begin{solution}{exo:g11:reaction-energy:6}
Rotos: $2 \times 436 + 498 = \qty{1370}{kJ}$; formados:
$4 \times 464 = \qty{1856}{kJ}$; $E_r \approx 1370 - 1856 =
\qty{-486}{kJ}$: \omterm{def:g11:reaction-energy:exothermic}{exotérmica}.
\end{solution}

\begin{solution}{exo:g11:reaction-energy:7}
El etanol \ce{CH3-CH2-O-H} tiene 5 \ce{C-H}, 1 \ce{C-C}, 1 \ce{C-O} y 1
\ce{O-H}. Rotos: $5 \times 415 + 345 + 350 + 464 + 3 \times 498 =
\qty{4728}{kJ}$. Formados: $4 \times 741 + 6 \times 464 = \qty{5748}{kJ}$.
$E_r \approx \qty{-1020}{kJ}$, alrededor de una cuarta parte menor, en
valor absoluto, que los \qty{1367.6}{kJ} medidos.
\end{solution}

\begin{solution}{exo:g11:reaction-energy:8}
Propano: $2219.2 / 0.0440 = \qty{5.04e4}{kJ/kg} = \qty{50.4}{MJ/kg}$.
Metano: $890.7 / 0.0160 = \qty{55.7}{MJ/kg}$. Los dos, como en el
diagrama.
\end{solution}

\begin{solution}{exo:g11:reaction-energy:9}
$Q = 1500 \times 4.18 \times 85 = \qty{5.33e5}{J} = \qty{533}{kJ}$;
$n = 533 / 890.7 = \qty{0.598}{mol}$, es decir, $0.598 \times 16.0 =
\qty{9.6}{g}$ de metano; si se pierde la mitad de la energía,
\qty{19}{g}.
\end{solution}

\begin{solution}{exo:g11:reaction-energy:10}
Hidrógeno > metano > propano > octano > etanol.
$48.0 / 142.9 = \qty{0.34}{kg}$ de hidrógeno.
\end{solution}

\begin{solution}{exo:g11:reaction-energy:11}
Rotos: $946 + 3 \times 436 = \qty{2254}{kJ}$; formados:
$6 \times 390 = \qty{2340}{kJ}$; $E_r \approx \qty{-86}{kJ}$: ligeramente
\omterm{def:g11:reaction-energy:exothermic}{exotérmica}.
\end{solution}

\begin{solution}{exo:g11:reaction-energy:12}
Agua: $250 \times 4.18 \times 20 = \qty{2.09e4}{J} = \qty{20.9}{kJ}$.
Etanol: $1.20 / 46.0 = \qty{0.0261}{mol}$, que podría liberar
$0.0261 \times 1367.6 = \qty{35.7}{kJ}$. \omterm{def:g10:synthesis-yield:yield}{Rendimiento}:
$20.9 / 35.7 \approx \qty{59}{\percent}$. Pérdidas: el calor que se lleva
el \omterm{def:g4:air-a-mixture-of-gases:air}{aire}, el calor que calienta la lata y el soporte, la \omterm{def:g8:combustion-and-fuels:complete}{combustión
incompleta} (hollín), algo de etanol que se \omterm{def:g3:separating-mixtures:evaporate}{evapora} sin \omterm{def:g4:what-a-fire-needs:burning}{arder}.
\end{solution}

\begin{solution}{exo:g11:reaction-energy:13}
\ce{2C8H18 + 25O2 -> 16CO2 + 18H2O}: 8 \omterm{def:g10:the-mole:amount}{moles} de dióxido de carbono,
\qty{352}{g}, por cada \qty{5470}{kJ}, así que $352 / 5.470 = \qty{64.4}{g}$ por megajulio. \ce{C3H8 + 5O2 -> 3CO2 + 4H2O}:
\qty{132}{g} por cada \qty{2219.2}{kJ}, así que \qty{59.5}{g} por
megajulio. El metano, con \qty{49.4}{g}, es el que menos libera.
\end{solution}

\begin{solution}{exo:g11:reaction-energy:14}
La tabla da \omterm{def:g11:reaction-energy:bond-energy}{energías de enlace} medias, mientras que los enlaces \ce{C=O}
del dióxido de carbono son más fuertes que la media; los valores medidos
corresponden al agua líquida, cuya condensación libera más energía que la
reacción en fase gaseosa; y el método trata cada enlace como
independiente de sus vecinos. La estimación sigue siendo útil: da el
signo correcto (\omterm{def:g11:reaction-energy:exothermic}{exotérmica}) y el orden de magnitud correcto sin ninguna
medida.
\end{solution}

\begin{solution}{exo:g11:reaction-energy:15}
Rotos: $4 \times 464 = \qty{1856}{kJ}$; formados: $2 \times 436 + 498 =
\qty{1370}{kJ}$; $E_r \approx \qty{+486}{kJ}$: \omterm{def:g11:reaction-energy:exothermic}{endotérmica}, así que hay
que aportar energía. Al quemar después el hidrógeno se recupera esa
energía: el hidrógeno almacena energía producida en otro lugar, no crea
ninguna.
\end{solution}

\begin{solution}{pb:g11:reaction-energy:1}
\textbf{1.} \ce{CH4 + 2O2 -> CO2 + 2H2O}.

\textbf{2.} \ce{C2H6O + 3O2 -> 2CO2 + 3H2O}.

\textbf{3.} \ce{2C8H18 + 25O2 -> 16CO2 + 18H2O}.

\textbf{4.} \ce{2H2 + O2 -> 2H2O}.

\textbf{5.} El hidrógeno.

\textbf{6.} \qty{16.0}{g/mol}, \qty{46.0}{g/mol}, \qty{114.0}{g/mol},
\qty{2.0}{g/mol}.

\textbf{7.} En \si{kJ} por \si{kg} y después en \si{MJ/kg}: metano,
$890.7 / 0.0160$, \qty{55.7}{MJ/kg}; etanol, $1367.6 / 0.0460$,
\qty{29.7}{MJ/kg}; octano, $5470 / 0.1140$, \qty{48.0}{MJ/kg};
hidrógeno, $285.8 / 0.0020$, \qty{142.9}{MJ/kg}.

\textbf{8.} Hidrógeno > metano > octano > etanol.

\textbf{9.} El hidrógeno es un gas muy ligero: un kilogramo ocupa un
volumen enorme a presión normal, así que hay que comprimirlo en pesados
depósitos de alta presión.

\textbf{10.} Metano: 4 \ce{C-H}; dioxígeno: 2 \ce{O=O} rotos. Dióxido de
carbono: 2 \ce{C=O}; agua: $2 \times 2 = 4$ \ce{O-H} formados.

\textbf{11.} $4 \times 415 + 2 \times 498 = \qty{2656}{kJ}$.

\textbf{12.} $2 \times 741 + 4 \times 464 = \qty{3338}{kJ}$.

\textbf{13.} $E_r \approx 2656 - 3338 = \qty{-682}{kJ}$, frente a los
\qty{-890.7}{kJ} medidos: alrededor de un \qty{23}{\percent} menor en
valor absoluto.

\textbf{14.} \omterm{def:g11:reaction-energy:bond-energy}{Energías de enlace} medias (el \ce{C=O} del dióxido de
carbono es más fuerte que la media); agua líquida en la medida, gases en
la estimación.

\textbf{15.} $1000 / 890.7 = \qty{1.123}{mol}$.

\textbf{16.} $1.123 \times 44.0 = \qty{49.4}{g}$.

\textbf{17.} Etanol: $1000 / 1367.6 = \qty{0.731}{mol}$, que dan
\qty{1.462}{mol} de dióxido de carbono, \qty{64.3}{g}. Octano:
$1000 / 5470 = \qty{0.183}{mol}$, que dan \qty{1.463}{mol},
\qty{64.4}{g}.

\textbf{18.} El metano: es el que tiene más \omterm{def:g7:atoms-and-molecules:atom}{átomos} de hidrógeno por
\omterm{def:g7:atoms-and-molecules:atom}{átomo} de carbono (4 frente a 2.25 del octano), y quemar hidrógeno da
energía sin dióxido de carbono.

\textbf{19.} De cómo se fabrica el hidrógeno: a partir del agua con
electricidad de fuentes renovables, casi no libera dióxido de carbono;
fabricado a partir del metano, libera el dióxido de carbono en la
fábrica en lugar de en el tubo de escape.

\textbf{20.} El metano libera unos \qty{49}{g} de dióxido de carbono por
megajulio.
\end{solution}
